% GENERATED FILE. Edit canonical lessons, then run npm run generate:books.
% canonical-source-hash: fnv1a64:40e9c71e6d885d6a
% canonical-lessons: FR-C408-contravention, FR-C408-radar, FR-C408-flasher, FR-C408-verbaliser, FR-C408-controler

\chapter{Contrôles sur la route}
\label{ch:fr-controles-sur-la-route}
\hlchaptermodality{408}

\begin{chapteropening}
\textbf{By the end of this chapter:} \emph{I can talk about road checks and fines: the fine, the speed camera, getting caught, writing someone up and checking.}

Complete the last lesson of chapter 408: I can talk about road checks and fines: the fine, the speed camera, getting caught, writing someone up and checking.
\end{chapteropening}

\section[la contravention]{la contravention — the fine (for a driving offence), the ticket}

\label{lesson:FR-C408-contravention}



\begin{quote}
Before the new one: say the French for the parallel parking space, then the French for the car pound.
\end{quote}

\subsection*{You'll want to know: la contravention}
\textbf{la contravention} — \textquotedblleft{}the fine (for a driving offence), the ticket\textquotedblright{}.

\begin{quote}
Il a reçu une contravention pour excès de vitesse.
\end{quote}

\begin{quote}
\emph{He got a ticket for speeding.}
\end{quote}

Informally \emph{un PV}, short for \emph{procès-verbal}.

\begin{cousinweb}
\emph{contravention} is a learned formation on Late Latin \emph{contrāvenīre} 'to oppose, to go against', from \emph{contrā} 'against' and \emph{venīre} 'to come' (Lewis \& Short).
\end{cousinweb}

\begin{grammarlens}[title={What you've built}]
One more word for telling, explaining and describing at B1.
\end{grammarlens}

\subsection*{Your turn}
\begin{itemize}
\raggedright
  \item \emph{Say it:} \emph{la contravention}
  \item \emph{Say it:} \emph{la contravention}, once more
  \item \emph{Say it:} a sentence of your own with \emph{la contravention}
  \item \emph{Recall:} say the French for the parallel parking space, then the French for the car pound, then say \emph{la contravention} again
\end{itemize}

\begin{tcolorbox}[breakable,colback=teal!4,colframe=teal!35!black,title={Before you move on}]
What does la contravention mean? (\textquotedblleft{}The fine (for a driving offence), the ticket\textquotedblright{}.) Say it once more.
\end{tcolorbox}
\section[le radar]{le radar — the speed camera; the radar}

\label{lesson:FR-C408-radar}



\begin{quote}
Before the new one: say the French for the car pound, then the French for the fine.
\end{quote}

\subsection*{You'll want to know: le radar}
\textbf{le radar} — \textquotedblleft{}the speed camera; the radar\textquotedblright{}.

\begin{quote}
Il y a un radar à l'entrée du village.
\end{quote}

\begin{quote}
\emph{There is a speed camera at the edge of the village.}
\end{quote}

\emph{Un radar fixe} stays in one place; \emph{un radar mobile} moves.

\begin{cousinweb}
\emph{radar} was borrowed from English \emph{radar}, made from the first letters of \emph{radio detection and ranging}. EtymDB lists the word without a source line.
\end{cousinweb}

\begin{grammarlens}[title={What you've built}]
One more word for telling, explaining and describing at B1.
\end{grammarlens}

\subsection*{Your turn}
\begin{itemize}
\raggedright
  \item \emph{Say it:} \emph{le radar}
  \item \emph{Say it:} \emph{le radar}, once more
  \item \emph{Say it:} a sentence of your own with \emph{le radar}
  \item \emph{Recall:} say the French for the car pound, then the French for the fine, then say \emph{le radar} again
\end{itemize}

\begin{tcolorbox}[breakable,colback=teal!4,colframe=teal!35!black,title={Before you move on}]
What does le radar mean? (\textquotedblleft{}The speed camera; the radar\textquotedblright{}.) Say it once more.
\end{tcolorbox}
\section[flasher]{flasher — to flash (of a speed camera), to catch speeding}

\label{lesson:FR-C408-flasher}



\begin{quote}
Before the new one: say the French for the fine, then the French for the speed camera.
\end{quote}

\subsection*{You'll want to know: flasher}
\textbf{flasher} — \textquotedblleft{}to flash (of a speed camera), to catch speeding\textquotedblright{}.

\begin{quote}
Je me suis fait flasher à quatre-vingt-dix.
\end{quote}

\begin{quote}
\emph{I got caught by a speed camera at ninety.}
\end{quote}

Informal; mostly \emph{se faire flasher}.

\begin{cousinweb}
\emph{flasher} is built on \emph{flash}, which was borrowed from English \emph{flash} (EtymDB): the camera flashes when you go too fast.
\end{cousinweb}

\begin{grammarlens}[title={What you've built}]
One more word for telling, explaining and describing at B1.
\end{grammarlens}

\subsection*{Your turn}
\begin{itemize}
\raggedright
  \item \emph{Say it:} \emph{flasher}
  \item \emph{Say it:} \emph{flasher}, once more
  \item \emph{Say it:} a sentence of your own with \emph{flasher}
  \item \emph{Recall:} say the French for the fine, then the French for the speed camera, then say \emph{flasher} again
\end{itemize}

\begin{tcolorbox}[breakable,colback=teal!4,colframe=teal!35!black,title={Before you move on}]
What does flasher mean? (\textquotedblleft{}To flash (of a speed camera), to catch speeding\textquotedblright{}.) Say it once more.
\end{tcolorbox}
\section[verbaliser]{verbaliser — to fine, to write up (a driver)}

\label{lesson:FR-C408-verbaliser}



\begin{quote}
Before the new one: say the French for the speed camera, then the French for to flash.
\end{quote}

\subsection*{You'll want to know: verbaliser}
\textbf{verbaliser} — \textquotedblleft{}to fine, to write up (a driver)\textquotedblright{}.

\begin{quote}
Le policier a verbalisé le conducteur.
\end{quote}

\begin{quote}
\emph{The police officer fined the driver.}
\end{quote}

Formal and official.

\begin{cousinweb}
\emph{verbaliser} is built on \emph{verbal}, as in \emph{procès-verbal} 'official written report', a learned word from Late Latin \emph{verbālis}, from \emph{verbum} 'word' (Lewis \& Short).
\end{cousinweb}

\begin{grammarlens}[title={What you've built}]
One more word for telling, explaining and describing at B1.
\end{grammarlens}

\subsection*{Your turn}
\begin{itemize}
\raggedright
  \item \emph{Say it:} \emph{verbaliser}
  \item \emph{Say it:} \emph{verbaliser}, once more
  \item \emph{Say it:} a sentence of your own with \emph{verbaliser}
  \item \emph{Recall:} say the French for the speed camera, then the French for to flash, then say \emph{verbaliser} again
\end{itemize}

\begin{tcolorbox}[breakable,colback=teal!4,colframe=teal!35!black,title={Before you move on}]
What does verbaliser mean? (\textquotedblleft{}To fine, to write up (a driver)\textquotedblright{}.) Say it once more.
\end{tcolorbox}
\section[contrôler]{contrôler — to inspect, to check; to control}

\label{lesson:FR-C408-controler}



\begin{quote}
Before the new one: say the French for to flash, then the French for to fine.
\end{quote}

\subsection*{You'll want to know: contrôler}
\textbf{contrôler} — \textquotedblleft{}to inspect, to check; to control\textquotedblright{}.

\begin{quote}
La police contrôle les papiers des automobilistes.
\end{quote}

\begin{quote}
\emph{The police check the drivers' papers.}
\end{quote}

\emph{Un contrôle} — a check.

\begin{cousinweb}
\emph{contrôler} comes from Old French \emph{contrerole} 'a second register kept to check the first', from Medieval Latin \emph{contrārotulum} (EtymDB), built on \emph{rotulus} 'roll', a small \emph{rota} 'wheel'.
\end{cousinweb}

\begin{grammarlens}[title={What you've built}]
One more word for telling, explaining and describing at B1.
\end{grammarlens}

\subsection*{Your turn}
\begin{itemize}
\raggedright
  \item \emph{Say it:} \emph{contrôler}
  \item \emph{Say it:} \emph{contrôler}, once more
  \item \emph{Say it:} a sentence of your own with \emph{contrôler}
  \item \emph{Recall:} say the French for to flash, then the French for to fine, then say \emph{contrôler} again
\end{itemize}

\begin{tcolorbox}[breakable,colback=teal!4,colframe=teal!35!black,title={Before you move on}]
What does contrôler mean? (\textquotedblleft{}To inspect, to check; to control\textquotedblright{}.) Say it once more.
\end{tcolorbox}
